\documentclass[10pt,a4paper]{article}
\usepackage[spanish,es-nodecimaldot]{babel}
\usepackage[utf8]{inputenc}
\usepackage[T1]{fontenc}
\usepackage{lmodern}
\usepackage[top=1.1cm,bottom=1.1cm,left=1.5cm,right=1.5cm]{geometry}
\usepackage{amsmath,amssymb}
\usepackage[table]{xcolor}
\usepackage{booktabs,array}
\usepackage{enumitem}
\usepackage[most]{tcolorbox}

\setlength{\parindent}{0pt}
\setlength{\parskip}{1.4pt}
\pagestyle{empty}

\AtBeginDocument{%
  \setlength{\abovedisplayskip}{3pt}%
  \setlength{\belowdisplayskip}{3pt}%
  \setlength{\abovedisplayshortskip}{2pt}%
  \setlength{\belowdisplayshortskip}{2pt}%
}

\begin{document}

\begin{center}
{\large\bfseries Tamaños muestrales distintos}
\end{center}
\vspace{-2pt}

cada fila tiene distintas cantidad de columnas

\begin{center}
\vspace{-2pt}
{\small
\begin{tabular}{l c c c c c c | c c c}
\toprule
Tratamiento & Obs. 1 & Obs. 2 & $\ldots$ & Obs. $j$ & $\ldots$ & Obs. $n_i$ & Tamaño ($n_i$) & Total ($T_{i.}$) & Media ($\bar{y}_{i.}$) \\
\midrule
Muestra 1 & $y_{11}$ & $y_{12}$ & $\ldots$ & $y_{1j}$ & $\ldots$ & $y_{1n_1}$ & $n_1$ & $T_{1.}$ & $\bar{y}_{1.}$ \\
Muestra 2 & $y_{21}$ & $y_{22}$ & $\ldots$ & $y_{2j}$ & --- & --- & $n_2$ & $T_{2.}$ & $\bar{y}_{2.}$ \\
$\ldots$ & $\ldots$ & $\ldots$ & $\ldots$ & $\ldots$ & $\ldots$ & $\ldots$ & $\ldots$ & $\ldots$ & $\ldots$ \\
Muestra $k$ & $y_{k1}$ & $y_{k2}$ & $\ldots$ & $y_{kn_k}$ & --- & --- & $n_k$ & $T_{k.}$ & $\bar{y}_{k.}$ \\
\midrule
Total Global & & & & & & & $N=\sum n_i$ & $T_{..}$ & $\bar{y}_{..}$ \\
\bottomrule
\end{tabular}
}
\end{center}

Se deben hacer los siguientes cambios:
\begin{enumerate}[leftmargin=20pt,itemsep=0pt,topsep=1pt,parsep=0pt]
  \item Cambiar n por $n_i$ en: SST, SS(Tr), C y T
  \item Cambiar k.n por N en: C (N es $=\sum n_i$)
\end{enumerate}
Se usa $n_i$ (tamaño de cada muestra) en lugar de n (tamaño de todas las muestras)
\[
\nu_1 := k-1 = \qquad\qquad \nu := N-1 = \qquad\qquad \nu_2 := \nu - \nu_1
\]

\[
C = \frac{T_{..}^2}{N} \quad \text{donde } N=\sum_{i=1}^{k} n_i
\]
\[
SST = \sum_{i=1}^{k}\sum_{j=1}^{n_i} y_{ij}^2 - C
\]
\[
SS(Tr) = \sum_{i=1}^{k} \frac{T_{i.}^2}{n_i} - C
\]
\[
SSE = SST - SS(Tr)
\]

\begin{center}
\vspace{-2pt}
{\small
\begin{tabular}{l c c c c}
\toprule
\textbf{Fuente de Variación} & \textbf{gl} & \textbf{Suma de Cuadrados} & \textbf{Media Cuadrada (MS)} & $F_{calc}$ \\
\midrule
Tratamientos & $k-1$ & $SS(Tr)$ & $MS(Tr)=\frac{SS(Tr)}{k-1}$ & $\frac{MS(Tr)}{MSE}$ \\
Error Aleatorio & $N-k$ & $SSE$ & $MSE=\frac{SSE}{N-k}$ & \\
\midrule
Total & $N-1$ & $SST$ & & \\
\bottomrule
\end{tabular}
}
\end{center}

\vspace{2pt}
\hrule
\vspace{4pt}

\begin{tcolorbox}[colback=blue!5, colframe=blue!45!black, title={\small\bfseries Enunciado del Problema}, boxsep=1pt, left=4pt, right=4pt, top=2pt, bottom=2pt]
{\small
El contenido de aflatoxina (en partes por millón, ppm) de algunas muestras de crema de maní de dos marcas comerciales dio los siguientes resultados no balanceados:

\vspace{2pt}
\begin{center}
\begin{tabular}{l c c c c c c c c | c c c}
\toprule
Marca & 1 & 2 & 3 & 4 & 5 & 6 & 7 & 8 & $n_i$ & Total ($T_{i.}$) & Media ($\bar{y}_{i.}$) \\
\midrule
Marca A & 0.5 & 0.0 & 3.2 & 1.4 & 0.0 & 1.0 & 8.6 & 2.9 & $n_1=8$ & 17.6 & 2.20 \\
Marca B & 4.7 & 6.2 & 0.0 & 10.5 & 2.1 & 0.8 & --- & --- & $n_2=6$ & 24.3 & 4.05 \\
\midrule
\multicolumn{9}{l}{Total} & $N=14$ & 41.9 & 2.99 \\
\bottomrule
\end{tabular}
\end{center}

1) Emplear ANOVA con $\alpha=0.05$ para probar si las marcas difieren en aflatoxina.

2) Probar la misma hipótesis usando la prueba $t$ de Student bimuestral y verificar la equivalencia $F=t^2$.
}
\end{tcolorbox}

\textbf{Resolución por ANOVA ($k=2, n_1=8, n_2=6, N=14$):}
\begin{enumerate}[leftmargin=20pt,itemsep=1pt,topsep=1pt,parsep=0pt]
  \item \textbf{Hipótesis:} $H_0: \mu_A=\mu_B$ vs $H_1: \mu_A \neq \mu_B$
  \item \textbf{Grados de Libertad y Valor Crítico:} $\nu_1=k-1=1$, $\nu_2=N-k=12$. Valor crítico $F_{0.05;1,12}=4.747$.
  \item \textbf{Cálculos:}
  \[
  C = \frac{41.9^2}{14} = \frac{1755.61}{14} = \mathbf{125.401}
  \]
  \[
  \sum\sum y_{ij}^2 = (0.5^2+\cdots+2.9^2)+(4.7^2+\cdots+0.8^2) = 93.91+177.74 = \mathbf{271.650}
  \]
  \[
  SST = 271.650-125.401 = \mathbf{146.249}
  \]
  \[
  SS(Tr) = \frac{17.6^2}{8}+\frac{24.3^2}{6}-125.401 = 38.72+98.415-125.401 = \mathbf{11.734}
  \]
  \[
  SSE = 146.249-11.734 = \mathbf{134.515}
  \]
  \item \textbf{Tabla ANOVA:}
  \begin{center}
  \vspace{-2pt}
  {\small
  \begin{tabular}{l c c c c}
  \toprule
  \textbf{Fuente} & \textbf{gl} & \textbf{Suma de Cuadrados} & \textbf{Media Cuadrada} & $F_{calc}$ \\
  \midrule
  Marcas (Trat.) & 1 & 11.734 & $MS(Tr)=11.734$ & \textbf{1.047} \\
  Error & 12 & 134.515 & $MSE=11.210$ & \\
  \midrule
  Total & 13 & 146.249 & & \\
  \bottomrule
  \end{tabular}
  }
  \end{center}
  \item \textbf{Decisión:} Como $F_{calc}=1.047 < F_{crit}=4.747$, \textbf{no se rechaza $H_0$}. No hay evidencia estadísticamente significativa para afirmar que las dos marcas difieran en su contenido medio de aflatoxina.
\end{enumerate}

\end{document}
